\Askhsh{22310_SOLUTION}{1}
\begin{ylh}{1}
    \item [ΛΥΣΗ]
    \begin{alist}
        \item 
        \begin{rlist}
            \item Επειδή η διάμετρος του ημικυκλίου είναι $AB = 8 \text{ cm}$, η ακτίνα του είναι
            \[ \rho = \frac{1}{2} AB = \frac{1}{2} \cdot 8 = 4 \text{ cm}, \]
            άρα το εμβαδόν του ημικυκλίου είναι
            \[ E = \frac{1}{2} \cdot \pi\rho^2 = \frac{1}{2} \cdot \pi \cdot 4^2 = 8\pi \text{ cm}^2. \]
            \item Το μήκος του ημικυκλίου είναι $L = \frac{1}{2} \cdot 2\pi\rho = \pi\rho = 4\pi \text{ cm}$.
        \end{rlist}
        \item 
        \begin{rlist}
            \item Το εμβαδόν του ορθογωνίου ΑΒΓΔ είναι $(\text{ΑΒΓΔ}) = AB \cdot ΑΔ = 8 \cdot ΑΔ \text{ cm}^2$.
            
            Επειδή το ημικύκλιο και το ορθογώνιο έχουν ίσα εμβαδά, θα έχουμε
            \[ (\text{ΑΒΓΔ}) = E \text{ ή } 8 \cdot ΑΔ = 8\pi \text{ ή } ΑΔ = \pi \text{ cm}. \]
            \item Το μήκος του ημικυκλίου από το ερώτημα (α ii) είναι $L = 4\pi \text{ cm}$.
            
            Στο (β i) βρήκαμε $ΑΔ = \pi \text{ cm}$ και επειδή το ΑΒΓΔ είναι ορθογώνιο, είναι $ΒΓ = ΑΔ = \pi \text{ cm}$ και $ΔΓ = AB = 8 \text{ cm}$.
            
            Η περίμετρος $P$ του σχήματος $(\Sigma)$ είναι
            \[ P = L + ΑΔ + ΔΓ + ΒΓ \quad \text{ή} \quad P = 4\pi + \pi + 8 + \pi = 6\pi + 8 \text{ cm}. \]
        \end{rlist}
    \end{alist}
\end{ylh}